\chapter{Multimodal AI for time series}
\label{chap:lecture_01}

These notes accompany the first half of the lecture, and follow the order of
the slides. They are not a complete treatment of multimodal learning. Each
section gives a short account of the main ideas, then points to the papers
the slides draw on and to good places to read further. The practical the next
day (ECG + EHR fusion) puts several of these ideas into code, and
Section~\ref{sec:lab} explains how.


\section{What makes a problem multimodal?}
\label{sec:intro}

A \emph{modality} is the way in which something happens or is experienced,
and a research problem is multimodal when it involves several such
modalities~\cite{baltrusaitis2019multimodal}. In machine learning the word is
used loosely for any distinct type of data: text, images, audio, waveforms,
tables, and so on. \emph{Multimodal machine learning} aims to build models that
can process and relate information from several of these at
once~\cite{baltrusaitis2019multimodal}.

What makes two data streams different modalities, rather than two channels of
one modality, is that they differ in structure, sampling and noise. The twelve
leads of an ECG are one modality: they are sampled together, at the same rate,
and measure the same physical process. An ECG and a patient's electronic
health record (EHR) are different modalities. The ECG is a dense 500\,Hz
waveform lasting ten seconds; the EHR is a sparse, irregular record of codes,
measurements and notes accumulated over years. Liang and
colleagues~\cite{liang2024foundations} describe multimodal data through three
principles: \emph{heterogeneity} (modalities look and behave differently),
\emph{connections} (they still carry related information about the same
thing), and \emph{interactions} (combining them can produce information that
neither has alone).

\paragraph{Why healthcare?}
Healthcare is multimodal by nature. A clinician rarely decides anything from
one source: a diagnosis brings together the history, an examination, blood
tests, imaging and physiological signals. The opportunities span discovery,
diagnosis, prognosis, treatment and prevention, and include remote monitoring
with wearable sensors, digital clinical trials and ``digital twins'' of
patients~\cite{acosta2022multimodal}. So do the challenges: data are scattered
across systems, incomplete, collected at different times, and subject to
privacy constraints.

\begin{furtherreading}
  \item Baltru\v{s}aitis, Ahuja \& Morency, \emph{Multimodal machine learning: a
        survey and taxonomy}~\cite{baltrusaitis2019multimodal}. The standard
        starting point, and the taxonomy used in
        Section~\ref{sec:challenges}.
  \item Acosta et al., \emph{Multimodal biomedical AI}~\cite{acosta2022multimodal}.
        A readable overview of the opportunities and obstacles in medicine.
  \item Liang, Zadeh \& Morency, \emph{Foundations and trends in multimodal machine
        learning}~\cite{liang2024foundations}. A more recent and more formal
        survey.
\end{furtherreading}


\section{A brief and subjective history}
\label{sec:history}

The idea that information from one sense changes how we perceive another is
old. A famous demonstration is the McGurk effect~\cite{mcgurk1976hearing}: when
the sound ``ba'' is dubbed over a video of someone saying ``ga'', many people
hear ``da''. Seeing the lips changes what we hear. It is fitting that one of the
earliest strands of multimodal machine learning was audio-visual speech
recognition, using lip movements to make speech recognition more robust.

A rough timeline:
\begin{itemize}
  \item \textbf{1990s and 2000s:} multimedia indexing and retrieval, finding
        images and videos by their content rather than by keywords alone.
  \item \textbf{2000s:} understanding human social interaction from recordings.
        The AMI Meeting Corpus~\cite{mccowan2005ami}, around 100 hours of meetings
        recorded with microphones, cameras and pens, is an example of the data
        this needed.
  \item \textbf{2010s:} deep learning arrives. Ngiam et al.~\cite{ngiam2011multimodal}
        showed that learning features from audio and video together could
        improve the features for each. Most of the decade's progress was on
        text and images: captioning, visual question answering, and later
        image--text pretraining (Section~\ref{sec:llms}).
  \item \textbf{Now:} almost anything is called a modality, and foundation models
        increasingly take several at once. Healthcare is one of the most
        active areas.
\end{itemize}

\begin{furtherreading}
  \item Baltru\v{s}aitis et al.~\cite{baltrusaitis2019multimodal}, Section~2, for a
        fuller history.
\end{furtherreading}


\section{Five challenges}
\label{sec:challenges}

Baltru\v{s}aitis et al.~\cite{baltrusaitis2019multimodal} organise multimodal
learning around five core challenges (Table~\ref{tab:challenges}). They overlap,
and other taxonomies exist~\cite{liang2024foundations}, but they are a useful
map. This lecture concentrates on \emph{representation} and \emph{fusion}, with
\emph{alignment} appearing in Section~\ref{sec:llms}.

\begin{table}[htbp]
  \centering
  \begin{tabular}{@{}lp{0.68\linewidth}@{}}
    \toprule
    Challenge      & The question it asks \\
    \midrule
    Representation & How should we represent and summarise data from each
                     modality, and from several together? \\
    Translation    & How can we map data from one modality to another, for
                     example an ECG to a written report? \\
    Alignment      & How do elements of one modality relate to elements of
                     another, for example which part of a report describes
                     which beat? \\
    Fusion         & How can we join information from several modalities to
                     make a prediction? \\
    Co-learning    & How can knowledge be transferred between modalities,
                     their representations and their models? \\
    \bottomrule
  \end{tabular}
  \caption{The five challenges of multimodal machine
           learning, after Baltru\v{s}aitis et al.~\cite{baltrusaitis2019multimodal}.}
  \label{tab:challenges}
\end{table}

\begin{furtherreading}
  \item Liang et al.~\cite{liang2024foundations} update the taxonomy to six
        challenges (representation, alignment, reasoning, generation,
        transference and quantification).
\end{furtherreading}


\section{Fusion}
\label{sec:fusion}

\subsection{Why fuse?}

There are two main reasons to combine modalities.
\begin{itemize}
  \item \textbf{Complementary information.} Each modality may carry information
        the others lack, so together they predict better than any one alone.
  \item \textbf{Robustness.} If one modality is missing or of poor quality, the
        others can make up for it.
\end{itemize}
Neither is guaranteed. Reviews of clinical work find that multimodal models
usually, but not always, beat unimodal ones, and that the gain depends on the
data and the task~\cite{kline2022multimodal,krones2025review}. Whether fusion
can help depends on how the information is shared:
some is \emph{redundant} (present in both modalities), some \emph{unique} to one,
and some \emph{synergistic} (only visible when the two are combined).
Liang et al.~\cite{liang2023quantifying} show how to estimate these quantities.

\subsection{Where to fuse: early, intermediate and late}

The most common way to describe fusion methods is by \emph{where} in the
pipeline the modalities meet (Figure~\ref{fig:fusion_paradigms}). Huang et
al.~\cite{huang2020fusion} use this taxonomy for imaging and EHR data, and
Krones et al.~\cite{krones2025review} for healthcare more broadly.

\begin{figure}[htbp]
  \centering
  \includegraphics[width=\linewidth]{generated/fusion_paradigms}
  \caption{Where the modalities meet. \emph{Early} fusion joins the input
           features and trains one model. \emph{Intermediate} fusion encodes
           each modality, combines the embeddings, and trains everything
           together. \emph{Late} fusion trains one model per modality and
           combines their predictions.}
  \label{fig:fusion_paradigms}
\end{figure}

\begin{itemize}
  \item \textbf{Early fusion} concatenates the input features (or features
        extracted from each modality) and fits a single model. It is simple and
        lets the model use interactions between modalities, but the inputs must
        be available together, and a high-dimensional modality can swamp a
        low-dimensional one.
  \item \textbf{Late fusion} trains a separate model on each modality and
        combines their outputs, by averaging, voting, or a small model trained
        on the outputs. It copes easily with a missing modality and can reuse
        existing models, but it cannot learn interactions between the
        modalities.
  \item \textbf{Intermediate} (or joint) fusion learns a representation of each
        modality and combines the representations inside one network, trained
        end to end. It is the most flexible, and needs the most data.
\end{itemize}
Real systems often mix these, for example fusing some modalities early and
adding others late~\cite{krones2025review}. Huang et al.~\cite{huang2020fusion}
also give practical guidelines for choosing between them.

\subsection{How to combine}

Within intermediate fusion, there are several ways to combine two embeddings
$\mathbf{h}_{\text{EHR}}$ and $\mathbf{h}_{\text{ECG}}$
(Figure~\ref{fig:fusion_operators}).

\begin{figure}[htbp]
  \centering
  \includegraphics[width=\linewidth]{generated/fusion_operators}
  \caption{Five ways to combine an EHR embedding (orange) with an ECG embedding
           (blue) into a fused representation (green).}
  \label{fig:fusion_operators}
\end{figure}

\begin{itemize}
  \item \textbf{Concatenation} $[\mathbf{h}_{\text{EHR}}, \mathbf{h}_{\text{ECG}}]$
        keeps everything and leaves the next layer to work out how to combine
        it.
  \item \textbf{Addition} $\mathbf{h}_{\text{EHR}} + \mathbf{h}_{\text{ECG}}$ keeps
        the size fixed, but assumes the two embeddings live in comparable
        spaces.
  \item \textbf{Gating} learns, for each patient, how much to trust each
        modality: $g \odot \mathbf{h}_{\text{EHR}} + (1 - g) \odot \mathbf{h}_{\text{ECG}}$,
        with $g \in [0, 1]$ computed from the inputs. Gated multimodal
        units~\cite{arevalo2017gated} are an early example; FiLM
        layers~\cite{perez2018film}, where one modality scales and shifts the
        features of another, are a close relative.
  \item \textbf{Multiplicative} combinations such as the outer product used by
        tensor fusion networks~\cite{zadeh2017tensor} model every pairwise
        interaction explicitly, at a high cost in parameters.
  \item \textbf{Attention} treats each modality (or each part of it) as a set of
        tokens and lets them attend to one another (Section~\ref{sec:transformers}).
  \item \textbf{A learnt shared space} trains the encoders so that matching
        inputs from the two modalities land close together, as in contrastive
        learning (Section~\ref{sec:clip}).
\end{itemize}

\subsection{Fusion in transformers}
\label{sec:transformers}

Transformers have become the default architecture for multimodal learning,
because once every modality is turned into a sequence of tokens, the same
attention machinery applies to all of them. Xu, Zhu \&
Clifton~\cite{xu2023multimodal} survey the design space. Broadly:
\begin{itemize}
  \item \textbf{early summation or concatenation} of the token sequences, fed to one
        transformer;
  \item \textbf{hierarchical} designs, which encode each modality separately and
        then fuse with further layers;
  \item \textbf{cross-attention}, where the queries come from one modality and
        the keys and values from another.
\end{itemize}
Cross-attention suits time series well: the Multimodal
Transformer~\cite{tsai2019multimodal} uses it to fuse streams recorded at
different sampling rates without aligning them first. Attention
bottlenecks~\cite{nagrani2021attention} restrict the exchange between
modalities to a few shared tokens, which reduces cost and often helps.

\subsection{Cautions}

Three findings are worth keeping in mind before building anything elaborate.
\begin{itemize}
  \item \textbf{Always compare with the best single modality.} A multimodal network
        trained end to end can do \emph{worse} than a unimodal one, because it
        overfits more easily and its modalities learn at different
        rates~\cite{wang2020what}. In jointly trained models the modalities can
        compete, so that only some are learnt well~\cite{huang2022modality}.
  \item \textbf{Missing modalities are the norm} in clinical data, and many models
        degrade sharply when one is absent~\cite{ma2022multimodal}.
  \item \textbf{More complex fusion is not always better.} A recent benchmark of
        EHR and chest X-ray fusion~\cite{yin2026when} found that fusion helps most
        when both modalities are present and complementary, and that the gains
        shrink quickly when data are missing.
\end{itemize}

\begin{furtherreading}
  \item Huang et al.~\cite{huang2020fusion} and Krones et
        al.~\cite{krones2025review}: fusion strategies in healthcare.
  \item Xu, Zhu \& Clifton~\cite{xu2023multimodal}: multimodal transformers.
  \item Wu et al.~\cite{wu2026deep}: a survey of learning with missing
        modalities.
  \item Clinical examples: HAIM~\cite{soenksen2022integrated} (embeddings from
        tabular data, time series, notes and images, combined for many tasks on
        MIMIC) and MedFuse~\cite{hayat2022medfuse} (clinical time series and chest
        X-rays, with missing images).
\end{furtherreading}


\section{EHRs as multimodal timelines}
\label{sec:ehr}

An EHR is itself multimodal and temporal: diagnosis and procedure codes,
laboratory results, medications, notes, imaging, signals and genomic data, each
recorded at irregular times over a patient's
life\todo{cite the source of the EHR timeline figure on slide 15}. A natural way
to model it is to treat the record as a sequence of \emph{events}, in the way a
language model treats text as a sequence of tokens.

\paragraph{Next-event prediction.} ETHOS~\cite{renc2024zero} tokenises each
patient's record into a timeline of medical events and trains a transformer to
predict the next event. Instead of a sequence of words, the model learns a
sequence of admissions, tests, results and diagnoses. Because it learns to
predict whole future trajectories, it can be used \emph{zero-shot}, without
labelled data or fine-tuning for each task, and can simulate how a patient's
course might unfold.

\paragraph{Time-to-event prediction.} MOTOR~\cite{steinberg2024motor} goes a
step further: instead of only the next event, it predicts the distribution of
the \emph{time} until each of many future events. This uses the temporal
structure in the data more fully, and suits clinical questions about how risk
changes over time. It was pretrained on 55 million patient records and adapts
to new prediction tasks with much less labelled data.

\paragraph{Bringing in more modalities.} Recent work extends these timelines
beyond structured codes. Kim et al.~\cite{kim2026computable} use language models
to extract structured clinical information from free-text notes, and combine it
with the structured record into computable patient journeys.
NOAH~\cite{susetzky2026noah} models the whole patient journey across the MIMIC
family of datasets, with modality-specific encoders for structured events,
text, chest X-rays, ECGs and echocardiograms, a time-aware transformer, and a
latent patient state.

\begin{furtherreading}
  \item Wornow et al.~\cite{wornow2023shaky}: a critical review of foundation models
        for EHRs, and how they should be evaluated.
  \item EHRSHOT~\cite{wornow2023ehrshot}: a benchmark for few-shot evaluation of
        EHR foundation models.
  \item Zhang et al.~\cite{zhang2023improving}: fusing irregularly sampled ICU time
        series with clinical notes.
\end{furtherreading}


\section{A worked example: ECGs and EHRs}
\label{sec:ecg_ehr}

ECGs are cheap, quick and recorded in huge numbers, and deep learning has shown
that they carry far more information than a clinician reads from them. A deep
network can predict one-year mortality from the 12-lead ECG
alone~\cite{raghunath2020prediction}, even among ECGs read as normal, and a
single ECG model can predict hundreds of diagnosis codes from the linked
EHR~\cite{strodthoff2024prospects}. The question for fusion is whether the EHR
adds anything the ECG does not already capture, and the other way round.

Khurshid et al.~\cite{khurshid2022ecg} give a clear clinical answer for
atrial fibrillation (AF). They trained a convolutional network on 12-lead ECGs
to predict five-year AF risk, and compared it with CHARGE-AF, an established
risk score built from clinical risk factors in the EHR. The combination of the
two, a simple form of late fusion, discriminated better than either alone
across three cohorts, including the UK Biobank. The ECG and the clinical risk
factors were complementary.

\subsection{In the practical}
\label{sec:lab}

The practical uses 1{,}000 ECGs from MIMIC-IV-ECG~\cite{gow2023mimicivecg},
linked to their machine-generated reports and to a small set of EHR fields from
MIMIC-IV, to predict mortality within one year of the ECG. It works through the
ideas in these notes:
\begin{enumerate}
  \item \textbf{Alignment.} Train a small CLIP-style model (Section~\ref{sec:clip})
        to align ECGs with their reports, using a 1D version of the ConvNeXt V2
        architecture~\cite{woo2023convnextv2}, and compare it with a pretrained
        ECG foundation model, ECGFounder\footnote{\url{https://github.com/PKUDigitalHealth/ECGFounder}}.
  \item \textbf{Unimodal baselines.} Predict mortality from the EHR alone and from
        the ECG alone.
  \item \textbf{Fusion.} Compare early, late and intermediate fusion, and the
        combination operators of Figure~\ref{fig:fusion_operators}.
\end{enumerate}

\begin{furtherreading}
  \item McKeen et al.~\cite{mckeen2025ecgfm}: ECG-FM, an open ECG foundation model
        evaluated on MIMIC-IV-ECG.
  \item Saporta et al.~\cite{saporta2024symile}: contrastive learning across more
        than two modalities, including ECGs, laboratory results and chest
        X-rays from MIMIC.
\end{furtherreading}


\section{Language models as multimodal backbones}
\label{sec:llms}

Large language models (LLMs) have become a common hub for multimodal systems.
Language is flexible enough to describe almost anything, and an LLM brings
reasoning and world knowledge for free. The question is how to get other
modalities in.

\subsection{Projecting into the language model}

The most widely copied recipe is LLaVA~\cite{liu2023visual}. Take an existing
encoder for the new modality (in LLaVA, a pretrained image encoder), add a small
learnt \emph{projector} that maps its embeddings into the LLM's token-embedding
space, and fine-tune on instructions. The LLM then reads the projected embeddings
as if they were words. Flamingo~\cite{alayrac2022flamingo} takes a different
route, inserting cross-attention layers into a frozen LLM so that its text tokens
can attend to the image features.

\subsection{Aligning representations: CLIP}
\label{sec:clip}

For a projector to work, the modality's embeddings must be compatible with the
language model's representations. CLIP~\cite{radford2021learning} learns such a
shared space directly, from pairs of images and captions. Given a batch of $N$
pairs, it encodes each image and each caption, normalises the embeddings
($\mathbf{e}_i$ and $\mathbf{t}_j$), and computes the similarities
$s_{ij} = \mathbf{e}_i \cdot \mathbf{t}_j / \tau$, where $\tau$ is a temperature.
Each image should match its own caption better than the others, so the loss is
a softmax cross-entropy over each row of the similarity matrix,
\begin{equation}
  L_{\text{image}\to\text{text}} = -\frac{1}{N} \sum_{i=1}^{N}
    \log \frac{\exp(s_{ii})}{\sum_{j=1}^{N} \exp(s_{ij})},
\end{equation}
and the same over each column. The CLIP loss is the average of the two
directions. This is the InfoNCE loss of van den Oord et
al.~\cite{vandenoord2018representation}. Trained on 400 million pairs from the
web, CLIP could classify images \emph{zero-shot}, by comparing an image with the
embeddings of captions such as ``a photo of a dog''.

The same idea transfers to medicine. MERL~\cite{liu2024zeroshot} aligns ECGs with
their reports, and can then classify ECGs zero-shot from text prompts.
ImageBind~\cite{girdhar2023imagebind} binds six modalities into one space, using
images as the hub.

\subsection{Medical examples}

\begin{itemize}
  \item \textbf{HeLM}~\cite{belyaeva2023multimodal} maps clinical data into an LLM's
        embedding space. Simple tabular data are written out as text; complex
        modalities, such as spirograms (a time series of lung function), pass
        through a learnt encoder. Adding the spirograms improved the
        prediction of asthma from UK Biobank data.
  \item \textbf{Med-PaLM M}~\cite{tu2024towards} is a single generalist model for
        clinical text, medical images and genomics, evaluated on 14 tasks.
\end{itemize}

\begin{furtherreading}
  \item Zhang et al.~\cite{zhang2024large}: a survey of LLMs for time series.
  \item Liu et al.~\cite{liu2025towards}: a survey of cross-modal modelling of time
        series and text.
  \item Wearable sensor data and language: Health-LLM~\cite{kim2024healthllm},
        PH-LLM~\cite{belyaeva2025personal} and SensorLM~\cite{zhang2025sensorlm}.
  \item ECG-specific multimodal LLMs: GEM~\cite{lan2025gem} and
        PULSE~\cite{liu2024teach}.
\end{furtherreading}


\section{The LLM context as unifier}
\label{sec:unifier}

Representation learning, alignment and fusion all need data and engineering
effort for each new modality. LLMs offer a shortcut, because they are famously
flexible about their input: anything that can be written down can go in the
prompt.

Ji et al.~\cite{ji2026portable} test how far this goes for clinical time series.
They use a frozen LLM to turn each patient's irregular ICU time series into a
short natural-language summary, then embed the summary with a frozen text
embedding model. The resulting representations perform competitively with
imputation, self-supervised learning and time series foundation models on
MIMIC-IV and two other ICU datasets, and lose less performance when moved to a
new hospital.

If text is a common currency, a further step is to let the LLM decide for
itself which data to fetch, which tools to run on them and how to combine the
results. That is the agentic approach, and the subject of the second half of the
lecture (Chapter~\ref{chap:lecture_02}).


\section{Summary}

\begin{itemize}
  \item Modalities differ in structure, sampling and noise, but carry related
        information. Healthcare data are multimodal by nature.
  \item Fusion can be early, intermediate or late, and embeddings can be combined
        by concatenation, addition, gating, attention or a learnt shared space.
        Always compare against the best single modality.
  \item EHRs can be modelled as multimodal timelines of events, with next-event
        and time-to-event objectives.
  \item LLMs act as multimodal backbones, either by projecting other modalities
        into their embedding space or by reading them as text.
\end{itemize}
